\subsection{Definition of a topological vector space}

DEFINITION 1. --- Given a topological division ring K (GT, III, § 6.7) and a set E such that E has

\(1^{\circ}\) the structure of a left vector space on \(\mathbf{K}\) ;

\(2^{\mathrm{o}}\) a topology compatible with the structure of the additive group of E (GT, III,

§ 1.1) and satisfying in addition the following axiom :

(EVT) the mapping \((\lambda, x) \mapsto \lambda x\) of \(\mathbf{K} \times \mathbf{E}\) in \(\mathbf{E}\) is continuous,

then E is called a left topological vector space over (or on) K.

It is equivalent to saying that E is a topological left K-module (GT, III, § 6.6).

A left vector space structure relative to K and a given topology on a set E, are said to be compatible if the topology and the additive group structure of E are compatible and if, in addition, the axiom (EVT) is valid. This is the same as saying that the two mappings \((x, y) \mapsto x + y\) and \((\lambda, x) \mapsto \lambda x\) of \(E \times E\) and of \(K \times E\) , respectively, in E are continuous, for then the mapping \(x \mapsto -x = (-1)x\) , is continuous and the topology of E is compatible with its additive group structure.

If \(\mathbf{E}\) is a left topological vector space over \(\mathbf{K}\) , we say that \(\mathbf{E}\) provided only with its vector space structure, underlies the topological vector space \(\mathbf{E}\) .

compatible with the vector space structure of E on R. For, let \(u_{0}\) be a point of E, let H be a compact subset of X and \(\varepsilon\) be an arbitrary strictly positive number. The real-valued function \(u_{0}\) is bounded in H; let \(a = \sup_{t \in H} |u_{0}(t)|\) ; if u is any point of E then for all \(t \in H\)

\[
\left| \lambda u (t) - \lambda_ {0} u _ {0} (t) \right| \leqslant | \lambda |. \left| u (t) - u _ {0} (t) \right| + a \left| \lambda - \lambda_ {0} \right|.
\]

Hence, if \(|\lambda - \lambda_0| \leqslant \varepsilon\) and \(|u(t) - u_0(t)| \leqslant \varepsilon\) for all \(t \in \mathbf{H}\) , then for \(t \in \mathbf{H}\) , \(|\lambda u(t) - \lambda_0 u_0(t)| \leqslant \varepsilon (\varepsilon + |\lambda_0| + a)\) , which shows that the axiom (EVT) is satisfied; similarly it can be verified that the compact convergence topology is compatible with the additive group structure of \(\mathbf{E}\) .

On the other hand, if X is not compact, the uniform convergence topology (in X) is not necessarily compatible with the vector space structure of E; for example if X = R and if \(u_{0}\) is an unbounded continuous function in R, then the mapping \(\lambda \mapsto \lambda u_{0}\) of R in E is not continuous in the uniform convergence topology on E.

\begin{enumerate}
\def\labelenumi{\arabic{enumi})}
\setcounter{enumi}{4}
\tightlist
\item
  Let E be a vector space of finite dimension n over a topological division ring K; there exists an isomorphism \(u:K_{s}^{n}\to E\) of vector K-spaces and moreover, if v is a second isomorphism of \(K_{s}^{n}\) on E, then we can write \(v=u\circ f\) , where f is an automorphism of the vector K-space \(K_{s}^{n}\) . Consider, on \(K_{s}^{n}\) , the product topology that is compatible with its vector space structure (Example 3); since every linear mapping of \(K_{s}^{n}\) in itself is continuous for this topology, every automorphism of the vector space \(K_{s}^{n}\) is bicontinuous. Hence, if we transfer the product topology of \(K_{s}^{n}\) to E, by means of any isomorphism whatever of \(K_{s}^{n}\) on E, the topology obtained on E is independent of the particular isomorphism used; we call it the canonical topology on E; we shall characterize it differently (I, § 1.3) when K is a non-discrete complete division ring with a valuation. Every linear mapping of E in a topological vector space over K is continuous for the canonical topology on E.
\end{enumerate}

In the same way as in def. 1, a right topological vector space over K, a topological division ring, can be defined; but every right vector space on K can be considered as a left vector space on the division ring \(K^{0}\) opposite to K (A, II, § 1.1) and the topology of K is compatible with the structure of the division ring \(K^{0}\) . For this reason we usually consider only left topological vector spaces; when we speak of «topological vector space» without qualification, it is to be understood that we refer to a left vector space.

If \(K'\) is a sub-division ring of K, and E a topological vector space over K, then it is clear that the topology of E is still compatible with the vector space structure of E relative to \(K'\) , obtained by restricting the field of scalars to \(K'\) ; we say that the topological vector space on \(K'\) , obtained by this procedure, underlies the topological vector space E on K.

In order that a topological vector space E be Hausdorff, it is necessary and sufficient that for all \(x \neq 0\) of E, there exists a neighbourhood of 0 not containing x (GT, III, § 1.2).

Consider a topology, on a vector space E over a topological division ring K, that is compatible with the additive group structure of E. Because of the identity

\[
\lambda x - \lambda_ {0} x _ {0} = (\lambda - \lambda_ {0}) x _ {0} + \lambda_ {0} (x - x _ {0}) + (\lambda - \lambda_ {0}) (x - x _ {0})
\]

axiom (EVT) is equivalent to the following system of three axioms.

(EVT') For all \(x_0 \in \mathrm{E}\) , the mapping \(\lambda \mapsto \lambda x_0\) is continuous at \(\lambda = 0\) .

(EVT \(_{II}\) ) For all \(\lambda_{0}\in K\) , the mapping \(x\mapsto\lambda_{0}x\) is continuous at x=0.

\((\mathrm{EVT}_{\mathrm{III}}^{\prime})\) The mapping \((\lambda, x) \mapsto \lambda x\) is continuous at \((0, 0)\) .

In particular :

PROPOSITION 1. --- For all \(\alpha \in \mathbf{K}\) and every point \(b \in \mathbf{E}\) , the mapping \(x \mapsto \alpha x + b\) of \(\mathbf{E}\) in itself is continuous. Further, if \(\alpha \neq 0\) , this mapping is a homeomorphism of \(\mathbf{E}\) on itself.

The second part of the proposition is a result of the fact that if \(\alpha \neq 0\) , then \(x \mapsto \alpha^{-1}x - \alpha^{-1}b\) is the inverse mapping of \(x \mapsto \alpha x + b\) .

COROLLARY. --- If A is an open (resp. closed) set in E, then \(\alpha A\) is open (resp. closed) in E for every \(\alpha \neq 0\) in K.

Let E and F be two topological vector spaces on the same topological division ring K. A bijection f of E on F is an isomorphism of the topological vector space E on the topological vector space F if and only if f is linear and bicontinuous. In particular, if \(\gamma \neq 0\) belongs to the centre of K, the homothety \(x \mapsto \gamma x\) is an automorphism of the topological vector space structure of E.

